% ==========================================================================
%  Chapter 1 — Introduction
% ==========================================================================
\chapter{Introduction}
\label{ch:introduction}

\epigraph{%
  The best way to understand a kernel is to run it where you can poke at it.%
}{Jeff Dike, 2000}

\section{What is User-Mode Linux?}
\label{sec:intro-what-is-uml}

User-Mode Linux (\UML{}) is an architecture port of the Linux kernel in which
the kernel itself executes as an ordinary userspace process on a host Linux
system.  Where a conventional kernel runs on bare metal---managing page tables,
servicing interrupts, and scheduling threads with the full authority of CPU ring
0---\UML{} performs all of these functions from ring 3, mediating access to host
resources through the host kernel's system call interface.

In the Linux source tree, \UML{} lives alongside x86, ARM, RISC-V, and the
other architecture ports under \umlpath{arch/um/}.  It shares the same
scheduler, the same VFS layer, the same memory management subsystem, and the
same networking stack as every other port.  The code that differs is precisely
the code that must differ: the low-level trap handling, context switching,
timer management, and memory mapping routines that, on hardware architectures,
speak directly to the CPU and chipset.  In \UML{}, these routines instead call
\syscall{clone}, \syscall{mmap}, \syscall{sigaction}, \syscall{ioctl}, and a
handful of other host system calls.

Jeff Dike created \UML{} in 1999--2000 and presented it at the USENIX Annual
Linux Showcase in Atlanta~\cite{dike2000}.  The original implementation used
the \texttt{ptrace} system call to intercept guest system calls: the \UML{}
kernel ran as a tracing process that caught every \texttt{SIGTRAP} delivered
when a guest userspace process invoked a system call.  This mechanism, while
conceptually elegant, imposed a heavy performance penalty---each guest syscall
required multiple context switches between the tracing kernel and the traced
guest---and the subsequent two decades of \UML{} development can be read as a
sustained effort to reduce that overhead.

The architecture is best understood through the Popek--Goldberg
framework~\cite{popek1974}.  A \emph{virtual machine monitor} (VMM)
must faithfully reproduce the behavior of the underlying machine for guest
software while retaining control over resource allocation.  Traditional VMMs
achieve this by exploiting hardware privilege rings: the guest kernel believes
it runs in ring 0, but the VMM intercepts sensitive instructions from a more
privileged mode.  \UML{} inverts this arrangement.  There is no privilege
separation \emph{at the hardware level} between the \UML{} kernel and its
guests; instead, the host kernel enforces separation through the ordinary
process isolation guarantees of Linux itself.  The \UML{} kernel and its guest
processes are all host-level processes.  What distinguishes them is that the
\UML{} kernel interposes on guest system calls before they can reach the
host---first via \texttt{ptrace}, later via \Seccomp{} filters, and now,
in the v2 redesign, optionally via \KVM{}.

\begin{figure}[ht]
\centering
\begin{tikzpicture}[
    layer/.style={
      draw=UMLGrid,
      fill=#1,
      minimum width=12cm,
      minimum height=1.2cm,
      rounded corners=4pt,
      font=\small
    },
    label/.style={
      font=\footnotesize\itshape,
      text=UMLMuted,
      anchor=west
    },
    arrow/.style={
      -{Stealth[length=5pt]},
      thick,
      color=UMLMuted
    }
  ]

  \node[layer=UMLBlue!8]  (host)    at (0,0)   {\textbf{Host Linux Kernel} (ring 0)};
  \node[layer=UMLTeal!8]  (backend) at (0,1.6)  {\textbf{Backend} (\texttt{ptrace} / \Seccomp{} / \KVMvTwo{})};
  \node[layer=UMLGreen!8] (uml)     at (0,3.2)  {\textbf{UML Kernel} (\umlpath{arch/um/})};
  \node[layer=KernelGold!8](guest)  at (0,4.8)  {\textbf{Guest Userspace}};

  \node[label] at (-6.4,0)   {Hardware privilege};
  \node[label] at (-6.4,1.6) {Syscall interposition};
  \node[label] at (-6.4,3.2) {Full Linux kernel};
  \node[label] at (-6.4,4.8) {Unmodified binaries};

  \draw[arrow] (guest.south) -- (uml.north)
    node[midway,right,font=\scriptsize] {guest syscall};
  \draw[arrow] (uml.south) -- (backend.north)
    node[midway,right,font=\scriptsize] {backend ops};
  \draw[arrow] (backend.south) -- (host.north)
    node[midway,right,font=\scriptsize] {host syscalls / \texttt{KVM\_RUN}};

\end{tikzpicture}
\caption[UML architectural layers]{%
  The four-layer \UML{} architecture.  Guest userspace programs issue system
  calls that are trapped by the \UML{} kernel.  The \UML{} kernel services
  these using standard Linux kernel subsystems, performing I/O through the
  backend and host kernel layers.
}
\label{fig:intro-layers}
\end{figure}


\section{Why UML Matters}
\label{sec:intro-why}

In a landscape crowded with virtualization and sandboxing technologies---\KVM{},
Xen~\cite{barham2003}, Firecracker~\cite{firecracker}, gVisor~\cite{gvisor},
containers, unikernels---it is reasonable to ask whether \UML{} still
has a role.  The answer is emphatically yes, precisely because \UML{} occupies
a unique and underserved niche in the design space.

\subsection{The real kernel, not a reimplementation}

gVisor's Sentry~\cite{gvisor} reimplements the Linux system call interface in
Go.  It is an impressive feat of engineering, but it is not the Linux kernel:
subtle semantic differences arise when the Sentry's interpretation of a system
call diverges from the kernel's.  The Linux Kernel Library
(\textsc{LKL})~\cite{lkl} compiles the kernel as a linkable library, but
strips away process scheduling and user/kernel separation.

\UML{}, by contrast, \emph{is} the Linux kernel.  Every syscall handler,
every scheduler decision, every lock acquisition follows the same code paths
as a kernel running on x86 hardware.  A bug found in \UML{} is a bug in
mainline Linux, and a fix validated in \UML{} is a fix for every architecture.
This fidelity is not accidental---it is a consequence of \UML{} being an
\texttt{arch/} port, sharing all architecture-independent code with the rest
of the tree.

\subsection{Full user/kernel separation}

Unlike \textsc{LKL}, which runs kernel code in the same address space and
privilege level as the application, \UML{} maintains genuine user/kernel
separation.  Guest userspace processes run in separate host processes with
restricted memory mappings.  A bug in a guest program cannot corrupt \UML{}
kernel data structures any more than a bug in a userspace program can corrupt
a bare-metal kernel.  This separation is essential for security research,
kernel fuzzing, and any application where the guest is untrusted.

\subsection{No hardware virtualization requirements}

\KVM{}~\cite{kivity2007} and Firecracker~\cite{agache2020} require hardware
virtualization extensions (Intel VT-x or AMD-V) and access to
\umlpath{/dev/kvm}.  In many environments---containerized CI/CD runners,
shared cloud instances, locked-down workstations---\texttt{/dev/kvm} is
unavailable.  \UML{} with the \texttt{ptrace} or \Seccomp{} backends has no
such requirement: it runs wherever a Linux userspace process can run.  This
makes it uniquely suited for environments where hardware-assisted
virtualization is either absent or administratively prohibited.

\subsection{Application domains}
\label{sec:intro-applications}

The combination of kernel fidelity, process isolation, and minimal host
requirements makes \UML{} valuable across a range of application domains:

\begin{description}[style=nextline,leftmargin=2em]
  \item[Kernel development and debugging.]
    A developer can compile the kernel, run it under \texttt{gdb}, set
    breakpoints in the scheduler, single-step through the page fault handler,
    and inspect kernel data structures---all from an ordinary terminal.
    No reboot, no serial console, no JTAG probe.

  \item[Kernel fuzzing.]
    Coverage-guided fuzzers like syzkaller~\cite{syzkaller} and
    kAFL~\cite{schumilo2017} benefit from fast boot times and deterministic
    execution.  \UML{} boots in under a second, can be snapshot/restored
    without hardware support, and when combined with \textsc{LKL}-style
    in-process execution for the fast path, can achieve millions of syscalls
    per second.

  \item[Network simulation.]
    Projects like Netkit~\cite{pizzonia2008} and Marionnet~\cite{marionnet}
    use \UML{} to run dozens or hundreds of virtual Linux machines, each with
    its own networking stack, connected by virtual switches and routers.
    Students and researchers can build and tear down complex network topologies
    entirely in userspace.

  \item[Security research.]
    \UML{} provides a complete kernel attack surface without requiring trust in
    hardware or a hypervisor.  Researchers can study kernel exploits, test
    sandboxing mechanisms, and develop mitigations in a fully controlled
    environment.

  \item[CI/CD and automated testing.]
    Kernel selftests, \texttt{kselftest} suites, and distribution build
    processes benefit from a lightweight kernel environment that can be
    instantiated in seconds inside a container or a GitHub Actions runner.

  \item[Education.]
    \UML{} turns the kernel into something a student can compile, run, modify,
    and observe without risk to the host system---lowering the barrier to
    kernel development from ``dedicated test machine'' to ``laptop terminal.''
\end{description}


\section{The Current Moment}
\label{sec:intro-moment}

This report is written at a moment of unusual convergence for \UML{}.  Five
developments, none individually decisive but collectively transformative,
have created a window of opportunity for the most significant redesign of
\UML{} since its creation.

\paragraph{1. The \Seccomp{} backend reaches mainline.}
The \Seccomp{} backend, which replaces the decades-old \texttt{ptrace}
mechanism with \texttt{SECCOMP\_IOCTL\_NOTIF\_SEND}-based syscall interception,
was merged for Linux 6.16.  This single change reduced guest syscall overhead
from approximately \SI{3}{\micro\second} (ptrace round-trip) to roughly
\SI{300}{\nano\second}, making \UML{} competitive with lightweight hypervisors
for I/O-intensive workloads.

\paragraph{2. SMP support lands.}
Historically, \UML{} was a uniprocessor architecture.  The addition of SMP
support---mapping \UML{} virtual CPUs to host threads with proper inter-processor
interrupt (IPI) emulation and TLB synchronization---allows \UML{} to exercise
the kernel's concurrency paths, which is precisely where the most dangerous
bugs hide.

\paragraph{3. gVisor's systrap validates swappable backends.}
Google's gVisor project demonstrated with its systrap platform~\cite{gvisor-systrap}
that a syscall-interposition system can support multiple backend mechanisms
behind a common interface.  The gVisor Sentry originally used \texttt{ptrace},
then \KVM{}, then systrap---each time improving performance while maintaining
the same guest-facing semantics.  This architectural precedent directly informed
the backend-ops abstraction in the \UML{} v2 redesign.

\paragraph{4. \textsc{LKL} shows what speed is possible.}
The Linux Kernel Library~\cite{lkl,purdila2010} compiles Linux as a
single-threaded library with no MMU or process separation.  It achieves
extraordinary fuzzing throughput---millions of operations per second---at
the cost of fidelity: there is no guest userspace, no scheduling, no realistic
attack surface.  \textsc{LKL} represents one extreme of the speed--fidelity
tradeoff.  The \UML{} v2 redesign aims to approach \textsc{LKL}'s speed while
retaining \UML{}'s complete kernel semantics.

\paragraph{5. The kernel community accepts AI-assisted development.}
The addition of the \texttt{coding-assistants.html} document to the kernel's
\umlpath{Documentation/process/} directory~\cite{coding-assistants} signals an
institutional willingness to engage with AI-assisted patch development.  The v2
redesign demonstrates what becomes possible when AI tools are used not to
generate mechanical patches, but to navigate architectural decisions across a
codebase of \num{30000000}+ lines of code.


\section{Report Structure}
\label{sec:intro-structure}

This report is organized as follows:

\begin{description}[style=nextline,leftmargin=2em,font=\bfseries]
  \item[Chapter~\ref{ch:introduction}: Introduction]
    (This chapter.)  Motivation, context, and scope.

  \item[Chapter 2: History]
    The evolution of \UML{} from Dike's original \texttt{ptrace}-based
    prototype through the TT/SKAS transition, the years of maintenance
    mode, and the recent resurgence of interest.

  \item[Chapter 3: Foundations]
    Background material on virtualization theory (Popek--Goldberg), Linux
    process primitives (\texttt{ptrace}, \Seccomp{}, \texttt{clone},
    \texttt{mmap}), and the \KVM{} API---the building blocks from which
    \UML{} is constructed.

  \item[Chapter 4: UML Architecture]
    The high-level architecture of \UML{}: how it maps kernel abstractions
    (virtual memory, interrupts, timers, I/O) onto host-level mechanisms.
    The backend-ops abstraction layer and its role in enabling pluggable
    backends.

  \item[Chapter 5: UML Implementation]
    A detailed walk through the \UML{} source tree, covering boot
    sequence, system call handling, memory management, signal delivery,
    SMP, and device I/O.

  \item[Chapter 6: The \KVMvTwo{} Backend]
    Design and implementation of the new \KVM{}-based backend.  How
    \texttt{KVM\_CREATE\_VM}, \texttt{KVM\_CREATE\_VCPU}, and
    \texttt{KVM\_RUN} are used to achieve approximately \SI{100}{\nano\second}
    syscall latency for guest-to-kernel transitions.

  \item[Chapter 7: Applications]
    Detailed case studies: kernel debugging under \texttt{gdb}, syzkaller
    integration, network simulation, CI/CD pipelines, and educational use.

  \item[Chapter 8: Related Work]
    Comparison with \KVM{}/QEMU, Xen, Firecracker, gVisor, \textsc{LKL},
    Dune~\cite{belay2012}, containers, and unikernels~\cite{madhavapeddy2013}.

  \item[Chapter 9: The v2 Redesign]
    The complete contribution of this work: backend-ops abstraction, \KVMvTwo{}
    backend, SMP, snapshot/restore, record/replay, comprehensive selftests,
    and the Rust \texttt{umlctl} launcher tool.  Design rationale, trade-offs,
    and lessons learned.

  \item[Chapter 10: Conclusion]
    Summary of contributions, open problems, and future directions.

  \item[Appendix A: Source Map]
    Annotated directory listing of the \UML{} source tree with the v2
    additions.

  \item[Appendix B: Benchmarks]
    Raw benchmark data and methodology for the performance measurements
    cited throughout the report.
\end{description}


\section{Contributions of This Work}
\label{sec:intro-contributions}

The v2 redesign of User-Mode Linux, developed by Michael Bommarito with
AI-assisted tooling, contributes the following:

\begin{enumerate}[leftmargin=2em]
  \item \textbf{Backend operations abstraction.}
    A \texttt{struct um\_backend\_ops} interface
    (defined in \umlpath{arch/um/include/shared/backend.h})
    that decouples the \UML{} kernel from any single syscall-interposition
    mechanism.  The abstraction defines a contract for backend lifecycle
    management, guest process creation, syscall interception, context switching,
    and memory mapping.  New backends can be added by implementing this
    interface without modifying core \UML{} code.

  \item \textbf{The \KVMvTwo{} backend.}
    A new backend that uses the \KVM{} API to run \UML{} guest code in
    hardware-accelerated virtual CPUs.  When the guest executes a system call,
    the resulting VM exit is handled directly by the \UML{} kernel---no
    \texttt{ptrace} round-trip, no \Seccomp{} notification.  Measured syscall
    latency is approximately \SI{100}{\nano\second}, a $3\times$ improvement
    over \Seccomp{} and a $30\times$ improvement over \texttt{ptrace}.

  \item \textbf{SMP support.}
    Multi-processor support for \UML{}, mapping virtual CPUs to host threads
    with proper IPI emulation, TLB shootdown, and memory barrier semantics.
    This enables \UML{} to exercise---and expose bugs in---the kernel's
    concurrent code paths.

  \item \textbf{Snapshot and restore.}
    Mechanisms for capturing and restoring the complete state of a running
    \UML{} instance, enabling fast restart for fuzzing workloads and
    reproducible bug investigation.

  \item \textbf{Record and replay.}
    Deterministic record/replay infrastructure that captures non-deterministic
    events (timer interrupts, I/O completions) during execution and replays
    them faithfully, enabling time-travel debugging of kernel code.

  \item \textbf{Comprehensive selftests.}
    A test suite under \umlpath{tools/testing/selftests/um/} covering backend
    lifecycle, syscall interception, SMP behavior, memory management, and
    performance regression detection.

  \item \textbf{The \texttt{umlctl} launcher.}
    A Rust-based command-line tool for building, configuring, launching, and
    managing \UML{} instances.  It handles kernel compilation, root filesystem
    preparation, backend selection, and instance lifecycle, reducing the barrier
    to entry from a page of shell commands to a single invocation.
\end{enumerate}

\noindent
These contributions are organized as seven series-level commits on
the \texttt{mjbommar/linux} tree (branch \texttt{next}), based on
\texttt{torvalds/master} at v7.1-rc5.  Each series is independently
reviewable and bisectable.

\bigskip

The remainder of this report proceeds from history and theory
(Chapters~2--3), through architecture and implementation (Chapters~4--6),
to applications and evaluation (Chapters~7--8), and concludes with
the complete v2 redesign narrative (Chapter~9) and future
directions (Chapter~10).
